```latex
\documentclass{article}
\usepackage{pathfinder-note}

\title{Costly Learning, Participation, and Balanced Transfers}

\begin{document}
\maketitle

\section{The pair in two sentences}
Q studies stochastic resource allocation with separable, strongly concave valuations and a quadratic payment rule intended to implement the welfare optimum with balanced transfers and individual rationality. P proposes charging LLM double-auction traders for generated tokens, but its internal labels ``Q'' and ``P'' refer to two other papers, so its auction is not an instance of the assigned Q's model. \ledger{1, 2, 8}

\section{Connexion pursued}
The useful connection is the cost of reaching an allocation. Q's individual-rationality argument concerns utility at an ideal Nash equilibrium, \(V_n(x_n^*)-t_n(s^*)\geq 0\); P separately accounts for the cost of decisions through \(W_{\lambda}=G-\lambda C\) and proposes charging even a trader's first decision. If Q's stochastic samples require paid, token-generating calls, that cost enters an agent's participation decision before the equilibrium benefit is realised. This is a conditional extension: Q does not say its samples are LLM calls, and P has not tested Q's mechanism. \ledger{3, 6, 21, 22, 24}

\section{What was found}
\begin{claim}
For one allocation with gross welfare \(G\), real metered cost \(\lambda C\), zero outside payoffs, and balanced transfers, individual rationality for every participant requires
\[
G\geq\lambda C.
\]
\end{claim}
Indeed, for payoffs \(u_n=g_n-\lambda c_n-t_n\), budget balance gives
\[
\sum_n u_n=\sum_n g_n-\lambda\sum_n c_n=G-\lambda C.
\]
If every \(u_n\geq0\), the sum must be nonnegative. This is an accounting condition, not an implementation or incentive theorem. Applied to P's symmetric auction schedule, whose gross-surplus ceiling is \(7.5\) per round, \(\lambda C>7.5\) rules out simultaneous budget balance and individual rationality for that round. \ledger{4, 11}

Q's full valuation class makes the participation issue concrete. With \(\mathcal X_n=[0,1]\) and \(V_n(x_n)=-\alpha x_n^2/2\) for every agent, the welfare optimum is zero and all of Q's stated concavity and normalization conditions hold. If entry requires a positive-cost sample before an agent can leave, total net welfare is negative; balanced transfers therefore cannot make every entrant individually rational. P's proposed withdrawal on the first scheduled call is charged, so it would not provide the costless pre-call choice needed to avoid this example. A pre-call opt-out, outside funding, or a restricted type class can remove this particular obstruction, without itself supplying a welfare guarantee. \ledger{15, 17, 21, 22, 24}

Even positive aggregate net value does not make Q's particular transfer cost-aware. For \(N=2\), \(\mathcal X_1=\mathcal X_2=[0,1]\), capacity \(c=2\), \(V_1(x)=-x^2/2\), and \(V_2(x)=x-x^2/2\), the optimum is \((0,1)\) with zero shadow price. Q's explicit payment rule admits the Nash profile \((x_1,x_2,p_1,p_2)=(0,1,0,0)\), where both transfers vanish. Agent 1 receives zero stage utility and would reject a mandatory cost of \(0.1\), although total net value remains \(0.5-0.1=0.4\). \ledger{5, 14}

There is also a finite-budget limit to the particular learning schedule in Q's Theorem 4. It sets \(Q_i=\lceil C_b/\eta^{2(i+1)}\rceil\), where \(0<\eta<1\). If each sample costs at least \(d>0\), all \(N\) agents sample at each iteration, and one allocation payoff is realised after \(I\) iterations, cumulative cost is at least
\[
dN\sum_{i=0}^{I-1}Q_i
\geq
\frac{dNC_b(\eta^{-2I}-1)}{1-\eta^2}.
\]
Compactness and continuity give a finite welfare bound \(M=\max_{x\in\mathcal X}\sum_n V_n(x_n)\). Budget balance and ex ante individual rationality can therefore hold only if
\[
I\leq
\frac{\log\!\left(1+M(1-\eta^2)/(dNC_b)\right)}
     {2\log(1/\eta)}.
\]
This bounds iterations under the stated cost model; it gives no accuracy guarantee at the stopping point. \ledger{6, 7, 12}

\section{Objections and scope}
Q's separable allocation theorem cannot be applied directly to P's buyer--seller matching. In a one-buyer, one-seller binary example, trade surplus \(h x_bx_s\) has cross-difference \(h>0\), while every additive \(V_b(x_b)+V_s(x_s)\) has zero cross-difference; binary participation also falls outside Q's convex allocation sets. \ledger{8, 16}

A fixed sampling charge that is independent of an agent's messages leaves Q's stage-game best responses unchanged after that charge has been incurred. A behavioural mechanism requires a choice about entry, withdrawal, sampling, or continuation. Q's simulations also state \(Q_i=\lceil0.96^{i+1}\rceil\), which equals one at every iteration and differs from its theorem's growing schedule; the plots cannot resolve the paid-sampling question. \ledger{7, 10}

\section{Outside sources and open question}
The targeted prior-work checks identified \href{https://academic.oup.com/ej/article/130/625/208/5593948}{\emph{Information Aggregation with Costly Reporting}}, \href{https://authors.library.caltech.edu/records/dxhp4-jvk64}{\emph{Optimal Data Acquisition for Statistical Estimation}}, \href{https://people.csail.mit.edu/silvio/Selected%20Scientific%20Papers/Mechanism%20Design/Delegating%20Computation%20On%20Costly%20Data.pdf}{\emph{Delegating Computation on Costly Data}}, and \href{https://arxiv.org/abs/1103.6280}{\emph{Bayesian Mechanism Design for Budget-Constrained Agents}} as adjacent work. The ledger does not establish novelty against them. \ledger{9, 13}

The material supports a research question, not a finished mechanism result: with endogenous, paid sampling and a pre-call entry choice, can an extension of Q achieve an explicit finite-budget approximate-Nash or welfare guarantee together with balanced transfers and ex ante individual rationality? The smallest next step is to specify who observes and pays each sample cost, when exit is available, and the finite-horizon payoff, then prove or refute those properties in that model. \ledger{16, 19, 24}

\end{document}
```